\begin{name}
	{\tenchude}{\tendethi}{LỚP TOÁN THẦY PHÁT}{\thoigian}
\end{name}
\setcounter{ex}{0}\setcounter{bt}{0}
\Opensolutionfile{ans}[ans/ans-2-TT-2-So-BacNinh-23]

\begin{ex}%[TT-SoBacNinh-2022-2023]%[Tran Tony, 12EX6-2023]%[2D1Y2-2]
	Cho hàm số $y=f(x)$ có bảng biến thiên như sau
	\begin{center}
		\begin{tikzpicture}
			\tkzTabInit[nocadre=false,lgt=1.2,espcl=2.5,deltacl=0.6]
			{$x$ /0.6, $y'$ /0.6, $y$ /2}
			{$-\infty$,$0$,$3$,$+\infty$}
			\tkzTabLine{ ,-,$0$,+,$0$,-, }
			\tkzTabVar{+/$+\infty$,-/$2$,+/$4$,-/$-\infty$}
		\end{tikzpicture}
	\end{center}
	Giá trị cực đại của hàm số $y=f(x)$ bằng
	\choice
	{$3$}
	{$0$}
	{$2$}
	{\True $4$}
	\loigiai{
		Quan sát bảng biến thiên, ta thấy giá trị cực đại của hàm số $y=f(x)$ bằng $4$.
	}
\end{ex}


\begin{ex}%[TT-SoBacNinh-2022-2023]%[Tran Tony, 12EX6-2023]%[2D1Y5-3]
	Cho hàm số $y=f(x)$ có bảng biến thiên như sau
	\begin{center}
		\begin{tikzpicture}
			\tkzTabInit[nocadre=false,lgt=1.2,espcl=2.5,deltacl=0.6]
			{$x$ /0.6, $y'$ /0.6, $y$ /2}
			{$-\infty$,$-1$,$0$,$2$,$+\infty$}
			\tkzTabLine{ ,+,$0$,-,$0$,+,$0$,-, }
			\tkzTabVar{-/$-\infty$,+/$5$,-/$1$,+/$3$,-/$-\infty$}
		\end{tikzpicture}
	\end{center}
	Phương trình $f(x)=m$ có bốn nghiệm thực phân biệt khi và chỉ khi
	\choice
	{$-1<m<2$}
	{$ 3<m<5$}
	{$ 1<m<5$}
	{\True $ 1<m<3$}
	\loigiai{
		Từ bảng biến thiên suy ra phương trình $f(x)=m$ có bốn nghiệm thực phân biệt khi và chỉ khi $1<m<3$.
	}
\end{ex}


\begin{ex}%[TT-SoBacNinh-2022-2023]%[Tran Tony, 12EX6-2023]%[2H2Y1-2]
	Cho một khối trụ có độ dài đường sinh là $l$ và bán kính của đường tròn đáy là $r$. Diện tích toàn phần của khối trụ là
	\choice
	{\True $S_{\text{tp}}=2\pi r\left(l+r\right)$}
	{$S_{\text{tp}}=\pi r\left(l+r\right)$}
	{$S_{\text{tp}}=2\pi r\left(l+2r\right)$}
	{$S_{\text{tp}}=\pi r\left(2l+r\right)$}
	\loigiai{
		Ta có diện tích toàn phần khối trụ $S_{\text{tp}}=2\pi lr+2\pi r^2=2\pi r\left(l+r\right)$.
	}
\end{ex}


\begin{ex}%[TT-SoBacNinh-2022-2023]%[Tran Tony, 12EX6-2023]%[2H1Y3-2]
	Thể tích khối lập phương có cạnh $a$ bằng
	\choice
	{$\dfrac{2a^3}{3}$}
	{$\dfrac{a^3}{3}$}
	{$ 2a^3$}
	{\True $a^3$}
	\loigiai{
		Thể tích khối lập phương có cạnh $a$ bằng $a^3$.
	}
\end{ex}


\begin{ex}%[TT-SoBacNinh-2022-2023]%[Tran Tony, 12EX6-2023]%[2D4Y1-1]
	Phần ảo của số phức $z=-3+2i$ bằng
	\choice
	{$-2$}
	{\True $2$}
	{$3$}
	{$-3$}
	\loigiai{
		Phần ảo của số phức $ z=-3+2i$ bằng $2$.
	}
\end{ex}


\begin{ex}%[TT-SoBacNinh-2022-2023]%[Tran Tony, 12EX6-2023]%[2D1Y1-2]
	\immini{
		Cho hàm số $ y=f(x)$ có đồ thị như hình vẽ. Khẳng định nào sau đây đúng?
		\choice
		{Hàm số $y=f(x)$ nghịch biến trên khoảng $(1;+\infty)$}
		{Hàm số $y=f(x)$ nghịch biến trên khoảng $(-\infty ;0)$}
		{Hàm số $y=f(x)$ đồng biến trên khoảng $(-2;1)$}
		{\True Hàm số $y=f(x)$ đồng biến trên khoảng $(0;1)$}
	}{
		\begin{tikzpicture}[scale=0.7, font=\footnotesize, line join=round, line cap=round, >=stealth]
			\tikzset{every node/.style={scale=0.9}}
			\draw[->] (-2,0)--(3,0) node[below] {$x$};
			\draw[->] (0,-3)--(0,2.7) node[left] {$y$};
			\fill (0,0) circle (1.4pt) ($(0,0)+(45:4mm)$) node{$O$};
			\fill (-1,0) circle (1.4pt) ($(-1,0)+(-135:4mm)$) node{$-1$};
			\fill (1,0) circle (1.4pt) ($(1,0)+(-90:4mm)$) node{$1$};
			\fill (2,0) circle (1.4pt) ($(2,0)+(-45:4mm)$) node{$2$};
			\fill (0,-2) circle (1.4pt) ($(0,-2)+(0:4mm)$) node{$-2$};
			\begin{scope}
				\clip (-2.1,-3.1) rectangle (3.1,2.9);
				\draw[samples=200,domain=-1.3:0,smooth,variable=\x] plot (\x,{4.64*(\x)^2+2.64*(\x)-2});
				\draw[samples=200,domain=0:1.31,smooth,variable=\x] plot (\x,{-1.98*(\x)^2+3.98*(\x)-2});
				\draw[samples=200,domain=1.31:2.9,smooth,variable=\x] plot (\x,{1.45*(\x)^2-4.53*(\x)+3.27});
			\end{scope}
		\end{tikzpicture}
	}
	\loigiai{
		Quan sát đồ thị ta thấy, trên khoảng $(0;1)$ dáng đồ thị tăng từ trái qua phải.\\
		Vậy hàm số $y=f(x)$ đồng biến trên khoảng $(0;1)$.
	}
\end{ex}

\begin{ex}%[TT-SoBacNinh-2022-2023]%[Tran Tony, 12EX6-2023]%[2D3Y1-1]
	Mệnh đề nào sau đây đúng?
	\choice
	{\True $\displaystyle\int{\cos}x\mathrm{\,d}x=\sin x+C$}
	{$\displaystyle\int{\cos}x\mathrm{\,d}x=\cos x+C$}
	{$\displaystyle\int{\cos}x\mathrm{\,d}x=-\sin x+C$}
	{$\displaystyle\int{\cos}x\mathrm{\,d}x=-\cos x+C$}
	\loigiai{
		Theo bảng nguyên hàm các hàm số cơ bản thì mệnh đề đúng là $\displaystyle\int{\cos}x\mathrm{\,d}x=\sin x+C$.
	}
\end{ex}

\begin{ex}%[TT-SoBacNinh-2022-2023]%[Tran Tony, 12EX6-2023]%[2H3Y2-4]
	Trong KG $Oxyz$, cho mặt phẳng $(\alpha)\colon x+y+z-3=0$. Điểm nào sau đây không thuộc mặt phẳng $(\alpha)$?
	\choice
	{$Q(1;-1;3)$}
	{$N(1;2;0)$}
	{$M(1;1;1)$}
	{\True $P(1;1;2)$}
	\loigiai{
		Thay lần lượt tọa độ các điểm vào phương trình mặt phẳng ta thấy\\
		Với $N(1;2;0)$ thì $1+2+0-3=0$ (luôn đúng). Do đó điểm $N(1;2;0)$ thuộc mặt phẳng $(\alpha)$.\\
		Với $P(1;1;2)$ thì $1+1+2-3=0$ (vô lí). Do đó điểm $P(1;1;2)$ không thuộc mặt phẳng $(\alpha)$.\\
		Với $Q(1;-1;3)$ thì $1-1+3-3=0$ (luôn đúng). Do đó điểm $Q(1;-1;3)$ thuộc mặt phẳng $(\alpha)$.\\
		Với $M(1;1;1)$ thì $1+1+1-3=0$ (luôn đúng). Do đó điểm $M(1;1;1)$ thuộc mặt phẳng $(\alpha)$.
	}
\end{ex}

\begin{ex}%[TT-SoBacNinh-2022-2023]%[Tran Tony, 12EX6-2023]%[2H3Y1-3]
	Trong KG $Oxyz$, cho mặt cầu $(S)\colon x^2+y^2+z^2-2x+4y-1=0$. Tâm của mặt cầu $(S)$ là điểm
	\choice
	{$I_2\left(1;1;1\right)$}
	{\True $I_4\left(1;-2;0\right)$}
	{$I_3\left(-2;4;-1\right)$}
	{$I_1\left(1;-2;-1\right)$}
	\loigiai{
		Ta có $(S)\colon x^2+y^2+z^2-2x+4y-1=0\Leftrightarrow{\left(x-1\right)^2}+\left(y+2\right)^2+z^2=6$.\\
		Suy ra tâm của mặt cầu $(S)$ là $I_4\left(1;-2;0\right)$.
	}
\end{ex}

\begin{ex}%[TT-SoBacNinh-2022-2023]%[Tran Tony, 12EX6-2023]%[2D1Y5-4]
	Đồ thị hàm số $y=-x^4+2x^2+3$ cắt trục tung tại điểm có tung độ bằng
	\choice
	{\True $3$}
	{$0$}
	{$1$}
	{$2$}
	\loigiai{
		Với	$x=0\Rightarrow y=3$. Vậy đồ thị hàm số cắt trục tung tại điểm có tung độ bằng $3$.
	}
\end{ex}

\begin{ex}%[TT-SoBacNinh-2022-2023]%[Tran Tony, 12EX6-2023]%[2D2B2-2]
	Trên khoảng $\left(0;+\infty\right)$, hàm số $y=x^{\frac{1}{3}}$ có đạo hàm là
	\choice
	{$y'=\dfrac{3}{4}{x^{\frac{4}{3}}}$}
	{$y'=\dfrac{1}{3}{x^{\frac{2}{3}}}$}
	{\True $y'=\dfrac{1}{3\sqrt[3]{x^2}}$}
	{$y'=\dfrac{1}{3\sqrt[3]{x}}$}
	\loigiai{
		Ta có $y=x^{\frac{1}{3}}\Rightarrow y'=\dfrac{1}{3\sqrt[3]{x^2}}$.
	}
\end{ex}

\begin{ex}%[TT-SoBacNinh-2022-2023]%[Tran Tony, 12EX6-2023]%[1D2B2-1]
	Số tập con gồm hai phần tử của tập hợp $A$ có $10$ phần tử là
	\choice
	{$10$}
	{\True $45$}
	{$20$}
	{$90$}
	\loigiai{
		Số tập con gồm hai phần tử của tập $A$ có $10$ phần tử là số tổ hợp chập $2$ của $10$ phần tử. Do đó có $\mathrm{C}_{10}^2=45$ tập con.
	}
\end{ex}

\begin{ex}%[TT-SoBacNinh-2022-2023]%[Tran Tony, 12EX6-2023]%[2D1B1-1]
	Cho hàm số $ y=f(x)$ xác định trên $\mathbb{R}$ và có đạo hàm $$f'(x)=12x^{2023}(x+1)(3-x),\forall x\in\mathbb{R}.$$ Hàm số đã cho đồng biến trên khoảng nào sau đây?
	\choice
	{$\left(-\infty ;0\right)$}
	{\True $\left(-\infty ;-1\right)$}
	{$\left(-1;3\right)$}
	{$\left(3;+\infty\right)$}
	\loigiai{
		Ta có $f'(x)=0\Leftrightarrow 12x^{2023}\left(x+1\right)\left(3-x\right)=0\Leftrightarrow\left[\begin{aligned}
			& x=0\\ 
			& x=-1\\ 
			& x=3.\\ 
		\end{aligned}\right.$\\
		Bảng xét dấu $f'(x)$
		\begin{center}
			\begin{tikzpicture}
				\tkzTabInit[lgt=1.2,espcl=3]
				{$x$ /0.6, $f'(x)$ /0.7}
				{$-\infty$,$-1$,$0$,$3$,$+\infty$}
				\tkzTabLine{ ,+,$0$,-,$0$,+,$0$, -, }
			\end{tikzpicture}
		\end{center}
		Từ bảng xét dấu suy ra, hàm số đã cho đồng biến trên các khoảng $\left(-\infty ;-1\right)$ và $(0;3)$.
	}
\end{ex}

\begin{ex}%[TT-SoBacNinh-2022-2023]%[Tran Tony, 12EX6-2023]%[2D2B6-1]
	Bất phương trình $\log\left(x-1\right)\le 2$ có tập nghiệm là
	\choice
	{$\left(0;101\right]$}
	{$\left(-\infty ;101\right]$}
	{\True $\left(1;101\right]$}
	{$\left(1;3\right]$}
	\loigiai{
		$\log\left(x-1\right)\le 2\Leftrightarrow\left\{\begin{aligned}
			& x-1>0\\ 
			& x-1\le{10^2}\\ 
		\end{aligned}\right.\Leftrightarrow\left\{\begin{aligned}
			& x>1\\ 
			& x\le 101\\ 
		\end{aligned}\right.\Leftrightarrow 1<x\le 101$.\\
		Vậy tập nghiệm của bất phương trình đã cho là $S=\left(1;101\right]$.
	}
\end{ex}

\begin{ex}%[TT-SoBacNinh-2022-2023]%[Tran Tony, 12EX6-2023]%[2H3B1-1]
	Trong KG $Oxyz$, cho ba điểm $A\left(2;2;0\right)$, $B\left(a;b;c\right)$, $ I\left(3;1;1\right)$. Điểm $ I$ là trung điểm của đoạn thẳng $ AB$. Khẳng định nào sau đây đúng?
	\choice
	{$ a+b+c=4$}
	{$ a+b+c=8$}
	{\True $ a+b+c=6$}
	{$ a+b+c=2$}
	\loigiai{
		Theo giả thiết $\heva{&2\cdot 3=a+2\\ &2\cdot 1=b+2\\&2\cdot 1=c}\Rightarrow a+b+c=6+2+2-2-2=6$.
		}
\end{ex}

\begin{ex}%[TT-SoBacNinh-2022-2023]%[Tran Tony, 12EX6-2023]%[2D2B4-5]
	Các nhà tâm lí học sử dụng mô hình hàm số để mô phỏng quá trình học tập của một học sinh như sau $f(x)=K\left(1-\dfrac{1}{e^{vx}}\right)$, trong đó $K$ là tổng số đơn vị kiến thức học sinh phải học, $v$ (kiến thức/ngày) là tốc độ tiếp thu của học sinh, $x$ (ngày) là thời gian học, $f(x)$ là số đơn vị kiến thức đã học được sau $x$ ngày. Giả sử một học sinh cần phải học 35 đơn vị kiến thức. Biết rằng tốc độ tiếp thu của học sinh này là $v=0{,}28$. Hỏi học sinh đó sẽ nhớ được bao nhiêu đơn vị kiến thức sau $7$ ngày (kết quả làm tròn đến hàng đơn vị)?
	\choice
	{\True $30$}
	{$26$}
	{$31$}
	{$21$}
	\loigiai{
		Đơn vị kiến thức học sinh đó nhớ được sau $7$ ngày là $f(7)=35\left(1-\dfrac{1}{\mathrm{e}^{7\cdot 0,28}}\right)\approx 30$.
	}
\end{ex}

\begin{ex}%[TT-SoBacNinh-2022-2023]%[Tran Tony, 12EX6-2023]%[2D4B1-1]
	Cho số phức $z$ có điểm biểu diễn trên mặt phẳng tọa độ là $A(3;-4)$. Giá trị của $ |z|$ bằng
	\choice
	{\True $5$}
	{$25$}
	{$10$}
	{$\sqrt{5}$}
	\loigiai{
		Số phức $z$ có điểm biểu diễn là $ A(3;-4)\Rightarrow z=3-4i\Rightarrow |z|=\sqrt{3^2+(-4)^2}=5$.
	}
\end{ex}

\begin{ex}%[TT-SoBacNinh-2022-2023]%[Tran Tony, 12EX6-2023]%[2D3B1-1]
	Hàm số nào sau đây là một nguyên hàm của hàm số $y=12x^5$?
	\choice
	{$y=12x^6+5$}
	{$y=12x^4$}
	{$y=60x^4$}
	{\True $y=2x^6+3$}
	\loigiai{
		Ta có $(2x^6+3)'=12x^5$. Vậy hàm số $y=2x^6+3$ là một nguyên hàm của hàm số $y=12x^5$.
	}
\end{ex}

\begin{ex}%[TT-SoBacNinh-2022-2023]%[Tran Tony, 12EX6-2023]%[2D4B4-1]
	Cho biết phương trình $z^2+az+b=0$ (với $a,b\in\mathbb{R}$) có nghiệm là $3-2i$. Giá trị của $a+b$ bằng
	\choice
	{\True $7$}
	{$-7$}
	{$19$}
	{$-19$}
	\loigiai{
		Do phương trình $z^2+az+b=0$ (với $a,b\in\mathbb{R}$) có nghiệm $3-2i$, khi đó 
		\allowdisplaybreaks
		\begin{eqnarray*}
			&&\left(3-2i\right)^2+a\left(3-2i\right)+b=0\\
			&\Leftrightarrow&3a+b+5+\left(-2a-12\right)i=0\\
			&\Leftrightarrow&\left\{\begin{aligned}
				& 3a+b+5=0\\ 
				&-2a-12=0\\ 
			\end{aligned}\right.\\
			&\Leftrightarrow&\left\{\begin{aligned}
				& a=-6\\ 
				& b=13.\\ 
			\end{aligned}\right.
		\end{eqnarray*}
		Vậy $a+b=7$.
	}
\end{ex}

\begin{ex}%[TT-SoBacNinh-2022-2023]%[Tran Tony, 12EX6-2023]%[2H1B3-2]
	Cho hình chóp $S.ABCD$ có đáy $ABCD$ là hình vuông cạnh $a$, $SA=3a$ và $SA$ vuông góc với mặt phẳng đáy. Thể tích khối chóp $S.ABCD$ bằng
	\choice
	{$3a^3$}
	{$9a^3$}
	{\True $a^3$}
	{$\dfrac{a^3}{3}$}
	\loigiai{
		Thể tích khối chóp $S.ABCD$ bằng $V_{S.ABCD}=\dfrac{1}{3}SA\cdot S_{ABCD}=a^3$.
	}
\end{ex}

\begin{ex}%[TT-SoBacNinh-2022-2023]%[Tran Tony, 12EX6-2023]%[2H3B3-3]
	Trong không gian $ Oxyz$, cho đường thẳng $d\colon\dfrac{x-1}{3}=\dfrac{y}{5}=\dfrac{z+2}{-2}$. Đường thẳng $d$ cắt mặt phẳng $\left(Oxy\right)$ tại điểm có hoành độ bằng
	\choice
	{$1$}
	{\True $-2$}
	{$-4$}
	{$-5$}
	\loigiai{
		Gọi $d\cap\left(Oxy\right)=M\left(x_0;{y_0};0\right)$.\\
		Thay tọa độ điểm $M$ vào $d\colon\dfrac{x-1}{3}=\dfrac{y}{5}=\dfrac{z+2}{-2}$ ta được $\dfrac{x_0-1}{3}=\dfrac{y_0}{5}=\dfrac{0+2}{-2}=-1\Rightarrow\left\{\begin{aligned}
			&{x_0}=-2\\ 
			&{y_0}=-5.\\ 
		\end{aligned}\right.$\\
		Vậy $M$ có hoành độ bằng $-2$.
	}
\end{ex}

\begin{ex}%[TT-SoBacNinh-2022-2023]%[Tran Tony, 12EX6-2023]%[2D2B4-2]
	Trên khoảng $\left(-\infty ;+\infty\right)$, hàm số $y=\ln\left(2x^2+1\right)$ có đạo hàm là
	\choice
	{$ y'=\dfrac{1}{2x^2+1}$}
	{\True $ y'=\dfrac{4x}{2x^2+1}$}
	{$ y'=4x\cdot\ln\left(2x^2+1\right)$}
	{$ y'=\dfrac{2x}{2x^2+1}$}
	\loigiai{
		Ta có $y'=\dfrac{\left(2x^2+1\right)'}{2x^2+1}=\dfrac{4x}{2x^2+1}$.
	}
\end{ex}

\begin{ex}%[TT-SoBacNinh-2022-2023]%[Tran Tony, 12EX6-2023]%[2D1B5-5]
	\immini{
		Cho hàm số $y=f(x)$ có đạo hàm trên $\mathbb{R}$. Đồ thị hàm số $y=f'(x)$ như hình vẽ. Khẳng định nào sau đây đúng?
		\choice
		{Hàm số $y=f(x)$ đạt cực đại tại điểm $x=-1$}
		{Hàm số $y=f(x)$ có đúng hai điểm cực trị}
		{\True Hàm số $y=f(x)$ đạt cực đại tại điểm $x=1$}
		{Hàm số $y=f(x)$ đạt cực đại tại điểm $x=4$}
	}{
		\begin{tikzpicture}[scale=0.7, font=\footnotesize, line join=round, line cap=round, >=stealth]
			\tikzset{every node/.style={scale=0.9}}
			\draw[->] (-2,0)--(4.8,0) node[below] {$x$};
			\draw[->] (0,-1.5)--(0,1.5) node[left] {$y$};
			\fill (0,0) circle (1.4pt) ($(0,0)+(-135:4mm)$) node{$O$};
			\foreach \x in {-1,1,4}
			\fill (\x,0) circle (1.4pt) ($(\x,0)+(-90:4mm)$) node{$\x$};
			\begin{scope}
				\clip (-1.8,-1.4) rectangle (5,1.5);
				\draw[samples=200,domain=-1.8:4.8,smooth,variable=\x] plot (\x,{0.15*((\x)^2-1)*((\x)-4)});
			\end{scope}
		\end{tikzpicture}
	}
	\loigiai{
		Quan sát đồ thị hàm số $y=f'(x)$, ta thấy $f'(1)=0$ và $f'(x)$ đổi dấu từ dương sang âm qua điểm $x=1$.\\
		Vậy hàm số $y=f(x)$ đạt cực đại tại điểm $x=1$.
	}
\end{ex}

\begin{ex}%[TT-SoBacNinh-2022-2023]%[Tran Tony, 12EX6-2023]%[2H3B3-2]
	Trong KG $Oxyz$, cho đường thẳng $d\colon\left\{\begin{aligned}
		& x=1+t\\ 
		& y=1+2t\\ 
		& z=4-3t\\ 
	\end{aligned}\right.$. Đường thẳng $d'$ đối xứng với $d$ qua mặt phẳng $\left(Oxy\right)$. PTTS của $d'$ là
	\choice
	{$\left\{\begin{aligned}
			& x=1+t'\\ 
			& y=1+2t'\\ 
			& z=-4-3t'\\ 
		\end{aligned}\right.$}
	{$\left\{\begin{aligned}
			& x=1+t'\\ 
			& y=1+2t'\\ 
			& z=0\\ 
		\end{aligned}\right.$}
	{\True $\left\{\begin{aligned}
			& x=1+t'\\ 
			& y=1+2t'\\ 
			& z=-4+3t'\\ 
		\end{aligned}\right.$}
	{$\left\{\begin{aligned}
			& x=-1-t'\\ 
			& y=-1-2t'\\ 
			& z=4-3t'\\ 
		\end{aligned}\right.$}
	\loigiai{
		Giả sử $d\cap\left(Oxy\right)=A\Rightarrow A\left(\dfrac{7}{3};\,\dfrac{11}{3};\,0\right)$ .\\
		Lấy điểm $M\left(1;\,1;\,4\right)\in d$ .\\
		Gọi $H$ là hình chiếu của $M$ xuống mặt phẳng $\left(Oxy\right)\Rightarrow H\left(1;\,1;\,0\right)$ .\\
		Gọi $M'$ đối xứng với $M$ qua mặt phẳng $\left(Oxy\right)\Rightarrow{M}'\left(1;\,1;\,-4\right)$ .\\
		Suy ra đường thẳng $d'$ đi qua hai điểm $A,\,M'$ .\\
		Ta có $\overrightarrow{M'A}=\left(\dfrac{4}{3};\,\dfrac{8}{3};\,4\right)=\dfrac{4}{3}\left(1;\,2;\,3\right)$ .\\
		Suy ra đường thẳng $d'$ đi qua điểm $M'\left(1;\,1;\,-4\right)$ , có véc-tơ chỉ phương $\overrightarrow{u}=\left(1;\,2;\,3\right)$ nên có phương trình là $\left\{\begin{aligned}
			& x=1+t'\\ 
			& y=1+2t'\\ 
			& z=-4+3t'.\\ 
		\end{aligned}\right.$
	}
\end{ex}

\begin{ex}%[TT-SoBacNinh-2022-2023]%[Tran Tony, 12EX6-2023]%[2D3B2-1]
	Nếu$\displaystyle\int\limits_1^2f(x)\mathrm{\,d}x=2$ thì $I=\displaystyle\int\limits_1^2\left[3f(x)-2\right]\mathrm{\,d}x$ bằng bao nhiêu?
	\choice
	{$I=3$}
	{\True $I=4$}
	{$I=2$}
	{$I=1$}
	\loigiai{
		Ta có $I=\displaystyle\int\limits_1^2\left[3f(x)-2\right]\mathrm{\,d}x=3\displaystyle\int\limits_1^2f(x)\mathrm{\,d}x-2\displaystyle\int\limits_1^2\mathrm{\,d}x=3\cdot 2-2x \bigg|_1^2=6-2=4$.
	}
\end{ex}

\begin{ex}%[TT-SoBacNinh-2022-2023]%[Tran Tony, 12EX6-2023]%[2D2B5-3]
	Tích các nghiệm của phương trình $6\log_4^2x-\log_4x^3+\dfrac{1}{5}=0$ bằng
	\choice
	{$\sqrt[3]{2}$}
	{$\dfrac{1}{30}$}
	{$4$}
	{\True $2$}
	\loigiai{
		Điều kiện $x>0.$\\
		Đặt $\log_4x=t$ thì phương trình trở thành $6t^2-3t+\dfrac{1}{5}=0$, do $\Delta=9-\dfrac{24}{5}>0$ nên phương trình luôn có $2$ nghiệm phân biệt, nên phương trình ban đầu luôn có $2$ nghiệm phân biệt.\\
		Ta có $x_1x_2=4^{t_1}{4^{t_2}}=4^{t_1+t_2}=4^{0{,}5}=2$.\\
		Vậy tích các nghiệm của phương trình ban đầu bằng $2$.
	}
\end{ex}

\begin{ex}%[TT-SoBacNinh-2022-2023]%[Tran Tony, 12EX6-2023]%[2H2B2-1]
	Nếu bán kính của một khối cầu tăng lên $2$ lần thì thể tích của khối cầu đó tăng lên bao nhiêu lần?
	\choice
	{$4$ lần}
	{\True $8$ lần}
	{$16$ lần}
	{$2$ lần}
	\loigiai{
		Ta có
		$$\begin{aligned}
			&{V_1}=\dfrac{4}{3}\pi R_1^3\\ 
			&{V_2}=\dfrac{4}{3}\pi R_2^3=\dfrac{4}{3}\pi{\left(2R_1\right)^3}=8\dfrac{4}{3}\pi R_1^3=8V_1.\\ 
		\end{aligned}$$
	}
\end{ex}

\begin{ex}%[TT-SoBacNinh-2022-2023]%[Tran Tony, 12EX6-2023]%[1H3B5-3]
	Cho hình chóp tứ giác đều $S.ABCD$ có $O$ là giao điểm của $AC$ và $BD$, $AB=SA=a$. Khoảng cách từ $O$ đến mặt phẳng $(SAB)$ bằng
	\choice
	{$\dfrac{a}{\sqrt{2}}$}
	{$\dfrac{a}{2}$}
	{$\dfrac{a\sqrt{3}}{2}$}
	{\True $\dfrac{a}{\sqrt{6}}$}
	\loigiai{
		\immini{
			Gọi $H$ là trung điểm $AB$, $K$ là hình chiếu vuông góc của $O$ lên $SH$.\\
			Ta có $AB\perp SO$, $AB\perp OH\Rightarrow AB\perp(SOH)\Rightarrow AB\perp OK$.\\
			Khi đó $AB\perp OK,OK\perp SH\Rightarrow OK\perp (SAD)$ hay $\mathrm{d}\left(O,(SAD)\right)=OK$.\\
			Ta có $SH=\dfrac{a\sqrt{3}}{2},\,OH=\dfrac{a}{2}\Rightarrow SO=\dfrac{a\sqrt{2}}{2}$.\\
			Vậy $\mathrm{d}\left(O,(SAD)\right)=OK=\dfrac{SO\cdot OH}{SH}=\dfrac{\dfrac{a\sqrt{2}}{2}\cdot\dfrac{a}{2}}{\dfrac{a\sqrt{3}}{2}}=\dfrac{a}{\sqrt{6}}$.
		}{
			\begin{tikzpicture}[scale=0.8, font=\footnotesize, line join=round, line cap=round, >=stealth]
				\coordinate (A) at (0,0);
				\coordinate (B) at (2,-2);
				\coordinate (D) at (5,0);
				\coordinate (C) at ($(B)+(D)-(A)$);
				\coordinate (O) at ($(A)!0.5!(C)$);
				\coordinate (S) at ($(O)+(0,4)$);
				\coordinate (H) at ($(A)!0.5!(B)$);
				\coordinate (K) at ($(S)!0.6!(H)$);
				\draw(S)--(A) (S)--(B) (H)--(S)--(C) (A)--(B)--(C);
				\draw[dashed,thin](A)--(C) (A)--(D) (C)--(D) (S)--(D) (S)--(O)--(H) (B)--(D) (O)--(K);
				\pic[draw,thin,angle radius=2mm] {right angle = S--O--D} 
				pic[draw,thin,angle radius=2mm] {right angle = S--K--O}
				pic[draw,thin,angle radius=2mm] {right angle = O--H--B};
				\foreach \i/\g in {S/90,A/180,B/-90,C/-90,D/0,O/-90,H/200,K/190}{\draw[fill](\i) circle (1.5pt) ($(\i)+(\g:3mm)$) node[scale=1]{$\i$};}
			\end{tikzpicture}
		}
	}
\end{ex}

\begin{ex}%[TT-SoBacNinh-2022-2023]%[Tran Tony, 12EX6-2023]%[2D4B4-1]
	Gọi $z_1$, $z_2$ là hai nghiệm của phương trình $z^2-4z+5=0$. Giá trị của $P=\left|z_1\right|^2+\left|z_2\right|^2$ là
	\choice
	{\True $P=10$}
	{$P=5$}
	{$P=6$}
	{$P=9$}
	\loigiai{
		Ta có $\Delta'=(-2)^2-1\cdot 5=-1<0$.\\
		Do đó phương trình đã cho có hai nghiệm phức phân biệt là $z_1=2+i$ và $z_1=2-i$.\\
		$P=\left|z_1\right|^2+\left|z_2\right|^2=\left| 2+i\right|^2+\left| 2-i\right|^2=\left(\sqrt{2^2+1^2}\right)^2+\left(\sqrt{2^2+\left(-1\right)^2}\right)^2=10$.
	}
\end{ex}

\begin{ex}%[TT-SoBacNinh-2022-2023]%[Tran Tony, 12EX6-2023]%[2D1B4-1]
	Giao điểm của đường tiệm cận đứng và đường tiệm cận ngang của đồ thị hàm số $y=\dfrac{3x-1}{x+2}$ là điểm nào sau đây?
	\choice
	{$Q(-1;2)$}
	{$N(3;-2)$}
	{\True $M(-2;3)$}
	{$P(2;-1)$}
	\loigiai{
		Đường tiệm cận đứng và tiệm cận ngang của đồ thị hàm số lần lượt là $x=-2$ và $y=3.$\\
		Nên giao điểm của $2$ đường tiệm cận là điểm $M(-2;3)$.
	}
\end{ex}

\begin{ex}%[TT-SoBacNinh-2022-2023]%[Tran Tony, 12EX6-2023]%[2D1B5-1]
	\immini{
		Hình vẽ bên là đồ thị của một trong bốn hàm số cho ở các phương án $A$, $B$, $C$, $D$. Hỏi đó là hàm số nào?
		\choice
		{$y=x^3-3x+1$}
		{\True $y=-x^3+3x+1$}
		{$y=-x^3+2x^2+1$}
		{$y=-x^3+3x^2+1$}
	}{
		\begin{tikzpicture}[scale=0.7, font=\footnotesize, line join=round, line cap=round, >=stealth]
			\tikzset{every node/.style={scale=0.9}}
			\draw[->] (-2.5,0)--(3,0) node[below] {$x$};
			\draw[->] (0,-1.5)--(0,3.5) node[left] {$y$};
			\fill (0,0) circle (1.4pt) ($(0,0)+(45:4mm)$) node{$O$};
			\fill (-1,0) circle (1.4pt) ($(-1,0)+(90:4mm)$) node{$-1$};
			\fill (1,0) circle (1.4pt) ($(1,0)+(-90:4mm)$) node{$1$};
			\fill (0,-1) circle (1.4pt) ($(0,-1)+(0:4mm)$) node{$-1$};
			\fill (0,1) circle (1.4pt) ($(0,1)+(180:4mm)$) node{$1$};
			\fill (0,3) circle (1.4pt) ($(0,3)+(180:4mm)$) node{$3$};
			\draw[dashed] (0,-1)--(-1,-1)--(-1,0);
			\draw[dashed] (0,3)--(1,3)--(1,0);
			\begin{scope}
				\clip (-2.3,-1.4) rectangle (2.3,3.5);
				\draw[samples=200,domain=-2.2:4.8,smooth,variable=\x] plot (\x,{-1*(\x)^3+3*(\x)+1});
			\end{scope}
		\end{tikzpicture}
	}
	\loigiai{
		Ta có hàm số dạng $y=ax^3+bx^2+cx+d$ có $y'=0$ có nghiệm $x=\pm 1$ và $a<0$.
	}
\end{ex}

\begin{ex}%[TT-SoBacNinh-2022-2023]%[Tran Tony, 12EX6-2023]%[1H3B3-3]
	Cho hình chóp $S.ABC$ có đáy là tam giác vuông cân tại $B$, $AB=a$, $SA=a\sqrt{2}$ và vuông góc với mặt phẳng đáy. Góc giữa đường thẳng $SC$ và mặt phẳng $(ABC)$ bằng
	\choice
	{\True $45^\circ $}
	{$30^\circ $}
	{$120^\circ $}
	{$60^\circ $}
	\loigiai{
		\immini{
		Ta có $ABC$ là tam giác vuông cân tại $B,\,AB=a$ nên $AC=a\sqrt{2}$.\\
		Tam giác $SAC$ có $SA=AC=a\sqrt{2}$ và $SA\perp AC$ (vì $SA$ vuông góc với mặt phẳng đáy $(ABC)$) nên $\Delta SAC$ vuông cân tại $A$.\\
		Từ đó suy ra $\left(SC,(ABC)\right)=\widehat{SCA}=45^{\circ}$.
		}{
			\begin{tikzpicture}[scale=0.9, font=\footnotesize, line join=round, line cap=round, >=stealth]
				\path
				(0,0) coordinate (A)
				(3,0) coordinate (C)
				(2,-1) coordinate (B)
				(0,3) coordinate (S)
				;
				\draw (A)--(B)--(C)--(S)--(A) (S)--(B);
				\draw[dashed] (A)--(C);
				\draw
				pic[draw, angle radius=2mm]{right angle=S--A--C}
				pic[draw, angle radius=2mm]{right angle=A--B--C}
				pic[draw, angle radius=5mm]{angle=S--C--A}
				;
				\foreach \x/\g in {A/180,B/220,C/70,S/180}
				\fill[black](\x) circle (1.3pt) ($(\x)+(\g:3mm)$) node{$\x$};
			\end{tikzpicture}
		}
	}
\end{ex}

\begin{ex}%[TT-SoBacNinh-2022-2023]%[Tran Tony, 12EX6-2023]%[1D2B5-2]
	Một hộp có $5$ viên bi màu đen, $4$ viên bi màu trắng. Chọn ngẫu nhiên $2$ viên bi. Xác suất chọn được $2$ bi cùng màu bằng
	\choice
	{$\dfrac{5}{9}$}
	{$\dfrac{1}{9}$}
	{$\dfrac{40}{9}$}
	{\True $\dfrac{4}{9}$}
	\loigiai{
		Số cách chọn $2$ bi là $n\left(\Omega\right)=\mathrm{C}_9^2$.\\
		Gọi $A$ là biến cố \lq\lq$2$ bi cùng màu\rq\rq.
		\begin{itemize}
			\item Chọn $2$ bi cùng màu đen có $\mathrm{C}_5^2$ cách.
			\item Chọn $2$ bi cùng màu trắng có $\mathrm{C}_4^2$ cách.
		\end{itemize}
		Suy ra $n(A)=\mathrm{C}_5^2+\mathrm{C}_4^2=16$.\\
		Vậy xác suất chọn được $2$ bi cùng màu là $\mathrm{P}(A)=\dfrac{n(A)}{n\left(\Omega\right)}=\dfrac{16}{36}=\dfrac{4}{9}$.
	}
\end{ex}

\begin{ex}%[TT-SoBacNinh-2022-2023]%[Tran Tony, 12EX6-2023]%[2D2B6-2]
	Bất phương trình $3^{\left|x\right|}<81$ có tập nghiệm là
	\choice
	{$\left[0;4\right]$}
	{$\left(0;4\right)$}
	{\True $\left(-4;4\right)$}
	{$\left(-\infty ;4\right)$}
	\loigiai{
		Ta có $3^{\left| x\right|}<81\Leftrightarrow{3^{\left| x\right|}}<3^4\Leftrightarrow\left| x\right|<4\Leftrightarrow-4<x<4$.	
	}
\end{ex}

\begin{ex}%[TT-SoBacNinh-2022-2023]%[Tran Tony, 12EX6-2023]%[2D3B2-1]
	Nếu $f(1)=12$ và $\displaystyle\int\limits_1^4f'(x)\mathrm{\,d}x=17$ thì giá trị của $f(4)$ bằng
	\choice
	{$ 9$}
	{$ 5$}
	{$ 19$}
	{\True $ 29$}
	\loigiai{
		Ta có $\displaystyle\int\limits_1^4f'(x)\text{d}x=f(4)-f(1)\Leftrightarrow 17=f(4)-12\Rightarrow f(4)=29$.
	}
\end{ex}

\begin{ex}%[TT-SoBacNinh-2022-2023]%[Tran Tony, 12EX6-2023]%[2H3B2-6]
	Trong KG $Oxyz$, cho hai mặt phẳng song song $(P)\colon 2x+y-2z-1=0$, $(Q)\colon 6x+3y-6z+15=0$. Khoảng cách giữa hai mặt phẳng $(P)$, $(Q)$ bằng
	\choice
	{\True $2$}
	{$\dfrac{16}{9}$}
	{$\dfrac{4}{3}$}
	{$\dfrac{16}{3}$}
	\loigiai{
		Ta có $(Q)\colon 6x+3y-6z+15=0\Leftrightarrow 2x+y-2z+5=0$. Lấy $M\left(0;1;0\right)\in (P)$.\\
		Vậy $\mathrm{d}\left((P),(Q)\right)=\mathrm{d}\left(M,(Q)\right)=\dfrac{\left| 1+5\right|}{\sqrt{1+4+4}}=2$.
	}
\end{ex}

\begin{ex}%[TT-SoBacNinh-2022-2023]%[Tran Tony, 12EX6-2023]%[1D3B3-3]
	Cho cấp số cộng $\left(u_n\right)$ có số hạng đầu $u_1=3$ và công sai $d=4$. Giá trị của $u_5$ bằng
	\choice
	{$-13$}
	{$768$}
	{$23$}
	{\True $19$}
	\loigiai{
		Ta có $u_5=u_1+4d=3+4\cdot 4=19$.
	}
\end{ex}

\begin{ex}%[TT-SoBacNinh-2022-2023]%[Tran Tony, 12EX6-2023]%[2D3B3-1]
	Gọi $(H)$ là hình phẳng giới hạn bởi các đường $ y=\sin x$; $Ox$; $x=0$; $x=\pi $. Diện tích của hình phẳng $(H)$ bằng
	\choice
	{$1$}
	{$\pi$}
	{\True $2$}
	{$2\pi$}
	\loigiai{
		Diện tích hình phẳng giới hạn bởi các đường $y=\sin x$, $Ox$, $x=0$, $x=\pi$ là $$S=\displaystyle\int\limits_0^{\pi}{\left|\sin x\right|}\mathrm{\,d}x=\displaystyle\int\limits_0^{\pi}{\operatorname{s}\text{in}x}\mathrm{\,d}x=-\cos x\bigg|_0^\pi=-\left(\cos\pi-\cos 0\right)=2.$$
	}
\end{ex}

\begin{ex}%[TT-SoBacNinh-2022-2023]%[Tran Tony, 12EX6-2023]%[2H2K1-2]
	Cho một hình trụ có bán kính đáy bằng $R$ và có chiều cao bằng $R\sqrt{3}$. Hai điểm $A$, $ B$ lần lượt nằm trên hai đường tròn đáy sao cho khoảng cách giữa đường thẳng $AB$ và trục của hình trụ bằng $\dfrac{R\sqrt{3}}{2}$. Góc giữa đường thẳng $AB$ và trục của hình trụ bằng
	\choice
	{$45^{^\circ}$}
	{$60^{^\circ}$}
	{$90^{^\circ}$}
	{\True $30^{^\circ}$}
	\loigiai{
		\immini{
			Kẻ đường sinh $AA'$, gọi $H$ là trung điểm của $A'B$.\\
			Ta có $\left\{\begin{matrix}
				O'H\perp A'B\\
				O'H\perp AA'\\
			\end{matrix}\right.\Rightarrow O'H\perp (ABA')$.\\
			$\mathrm{d}(AB;OO')=\mathrm{d}\left(OO';\left(ABA'\right)\right)=\mathrm{d}\left(O';\left(ABA'\right)\right)=O'H=\dfrac{R\sqrt{3}}{2}$ .\\
			$\Rightarrow A'B=2\sqrt{R^2-\left(\dfrac{R\sqrt{3}}{2}\right)^2}=R$.
		}{
			\begin{tikzpicture}[line join=round, line cap=round,thick]
				\def\a{2}
				\def\b{0.7}
				\def\h{3}
				\coordinate (M) at (0,0);
				\coordinate (N) at ($(M)+(2*\a,0)$);
				\coordinate (O) at ($(M)!0.5!(N)$);
				\coordinate (M') at ($(M)-(0,\h)$);
				\coordinate (N') at ($(N)-(0,\h)$);
				\coordinate (O') at ($(M')!0.5!(N')$);
				\draw(M)--(M') (N)--(N');
				\draw[dashed,thin](O)--(O');
				\draw[dashed,thin] (M') arc (180:0:\a cm and \b cm);
				\draw (O) ellipse (\a cm and \b cm)	(M') arc (-180:0:\a cm and \b cm);
				\path (N) arc (0:-76:{\a} and {\b})	coordinate (A)
				(N') arc (0:-76:{\a} and {\b})	coordinate (A')
				(N') arc (0:40:{\a} and {\b})	coordinate (B)
				($(A')!1/2!(B)$) coordinate (H)
				;
				\pic[draw,thin,angle radius=1.7mm] {right angle = O'--H--B}
				pic[draw,thin,angle radius=1.7mm] {right angle = A--A'--H};
				\draw (A)--(A');
				\draw[dashed] (A')--(O')--(B)--(A') (A)--(B) (O')--(H);
				\foreach \i/\g in {O/90,O'/-90,A/-130,A'/-120,B/50,H/-40}{\draw[fill](\i) circle (1.3pt) ($(\i)+(\g:2.6mm)$) node[scale=0.7]{$\i$};}
			\end{tikzpicture}
		}
		\noindent
		Góc giữa đường thẳng $AB$ và trục của hình trụ là góc giữa $AB$ và $AA'$, là góc $\widehat{BAA'}$\\
		Xét tam giác $ABA'$ vuông tại $A'$, có $\tan \widehat{BAA'}=\dfrac{A'B}{AA'}=\dfrac{R}{R\sqrt{3}}=\dfrac{1}{\sqrt{3}}\Rightarrow\widehat{BAA'}=30^0$.
	}
\end{ex}

\begin{ex}%[TT-SoBacNinh-2022-2023]%[Tran Tony, 12EX6-2023]%[2H3K2-8]
	Trong không gian $ Oxyz$, cho các điểm $ A(1;1;1)$, $B(-1;2;0)$, $C(3;-1;2)$ và mặt phẳng $(\alpha)\colon 2x-y+2z+7=0$. Điểm $M$ chạy tùy ý trên $(\alpha)$. Gọi $m$ là giá trị nhỏ nhất của biểu thức $\left| 3\overrightarrow{MA}+5\overrightarrow{MB}-7\overrightarrow{MC}\right|$. Khẳng định nào sau đây đúng?
	\choice
	{$m\in (28;47)$}
	{\True $m\in (24;28)$}
	{$ m\in (10;20)$}
	{$ m\in (20;24)$}
	\loigiai{
		Gọi $I$ là điểm thỏa mãn $3\overrightarrow{IA}+5\overrightarrow{IB}-7\overrightarrow{IC}=\overrightarrow{0}\Leftrightarrow I=\dfrac{3A+5B-7C}{3+5-7}=\left(-23;20;-9\right)$.\\
		Ta có $\mathrm{d}\left(I;(\alpha)\right)=\dfrac{\left|2\cdot\left(-23\right)-20+2\cdot\left(-9\right)+7\right|}{\sqrt{2^2+\left(-1\right)^2+2^2}}=\dfrac{77}{3}$.\\
		Xét $\left| 3\overrightarrow{MA}+5\overrightarrow{MB}-7\overrightarrow{MC}\right|=\left|\overrightarrow{MI}+\left(3\overrightarrow{IA}+5\overrightarrow{IB}-7\overrightarrow{IC}\right)\right|=MI\ge \mathrm{d}\left(I;(\alpha)\right)=\dfrac{77}{3}$.
	}
\end{ex}

\begin{ex}%[TT-SoBacNinh-2022-2023]%[Tran Tony, 12EX6-2023]%[2D3K2-2]
	Cho hàm số $ f(x)$ liên tục trên $\mathbb{R}$. Biết $\displaystyle\int_1^{\mathrm{e}^3}{\dfrac{f(\ln x)}{x}}\mathrm{\,d}x=7$,$\displaystyle\int_0^{\tfrac{\pi}{2}}{f}(\cos x)\sin x\mathrm{\,d}x=3$. Giá trị của $\displaystyle\int_1^3\left[f(x)+2x\right]\mathrm{\,d}x$ bằng
	\choice
	{$10$}
	{$15$}
	{$-10$}
	{\True $12$}
	\loigiai{
		Ta có $\displaystyle\int\limits_1^{\mathrm{e}^3}{\dfrac{f\left(\ln x\right)}{x}\mathrm{\,d}x}=\displaystyle\int\limits_1^{\mathrm{e}^3}{f\left(\ln x\right)\mathrm{\,d}\left(\ln x\right)}=\displaystyle\int\limits_0^3f(t)\mathrm{\,d}t=\displaystyle\int\limits_0^3f(x)\mathrm{\,d}x=7$.\\
		Ta có $\displaystyle\int\limits_0^{\tfrac{\pi}{2}}{f\left(\cos x\right)\sin x\mathrm{\,d}x}=-\displaystyle\int\limits_0^{\tfrac{\pi}{2}}{f\left(\cos x\right)\mathrm{\,d}\left(\cos x\right)}=-\displaystyle\int\limits_1^0f(u)\mathrm{\,d}u=\displaystyle\int\limits_0^1f(x)\mathrm{\,d}x=3$.\\
		Khi đó $\displaystyle\int\limits_1^3\left[f(x)+2x\right]\mathrm{\,d}x=\displaystyle\int\limits_0^3f(x)\mathrm{\,d}x-\displaystyle\int\limits_0^1f(x)\mathrm{\,d}x+\displaystyle\int\limits_1^32x\mathrm{\,d}x=7-3+8=12$.
	}
\end{ex}

\begin{ex}%[TT-SoBacNinh-2022-2023]%[Tran Tony, 12EX6-2023]%[2D1K1-3]
	Có bao nhiêu giá trị nguyên của tham số $m\in\left[-2022;2022\right]$ để hàm số $$y=-x^4+2m^2x^2+m^3$$ nghịch biến trên khoảng $\left(-4;0\right)?$
	\choice
	{$4036$}
	{$2019$}
	{$4045$}
	{\True $4038$}
	\loigiai{
		Có $y'=-4x^3+4m^2x=4x\left(m^2-x^2\right)=0\Leftrightarrow\hoac{&x=0\\&x=\pm m.}$\\
		Nếu $m=0$ ta có bảng xét dấu của $y'$:
		\begin{center}
			\begin{tikzpicture}
				\tkzTabInit[nocadre=false,lgt=1.2,espcl=2.5,deltacl=0.6]
				{$x$ /0.7, $y'$ /0.7}
				{$-\infty$,$0$,$+\infty$}
				\tkzTabLine{ ,+,$0$,-, }
			\end{tikzpicture}
		\end{center}
		Nên hàm số đồng biến trên khoảng $(-4;0)$ (không thỏa mãn).\\
		Nếu $m\ne 0$ , $y'$ có $3$ nghiệm phân biệt và có xét dấu như sau:
		\begin{center}
			\begin{tikzpicture}
				\tkzTabInit[nocadre=false,lgt=1.2,espcl=2.5,deltacl=0.6]
				{$x$ /0.7, $y'$ /0.7}
				{$-\infty$,$-|m|$,$0$,$|m|$,$+\infty$}
				\tkzTabLine{ ,+,$0$,-,$0$,+,$0$,- }
			\end{tikzpicture}
		\end{center}
		Để hàm sô nghịch biến trên khoảng $\left(-4;0\right)$ thì $-\left| m\right|\le-4\Leftrightarrow \hoac{&m\ge 4\\&m\le-4.}$\\
		Kết hợp giả thiết $m\in\left[-2022;2022\right]$ ta được $m\in\left[-2022;-4\right]\cup\left[4;2022\right]$.\\
		Số giá trị nguyên $m$ thỏa mãn đề bài là $4038$.
	}
\end{ex}

\begin{ex}%[TT-SoBacNinh-2022-2023]%[Tran Tony, 12EX6-2023]%[2H1K3-2]
	Cho hình lăng trụ tam giác đều $ABC.A'{B}'{C}'$ có cạnh đáy bằng $a$. Biết khoảng cách giữa hai đường thẳng $AB$ và $A'C$ bằng $\dfrac{a\sqrt{15}}{5}$. Thể tích $V$ của khối lăng trụ $ABC.A'{B}'{C}'$ là
	\choice
	{$ V=\dfrac{3a^3}{2}$}
	{$ V=\dfrac{3\sqrt{3}{a^3}}{8}$}
	{\True $ V=\dfrac{3a^3}{4}$}
	{$ V=\dfrac{3a^3}{8}$}
	\loigiai{
		\immini{
			Gọi $I$, $I'$ lần lượt là trung điểm của $AB$ và $A'B'$ suy ra $IC=\dfrac{a\sqrt{3}}{2}$.\\
			Gọi $H$ là hình chiếu của $I$ lên cạnh $I'C$ suy ra $\text{d}\left(I,(A'B'C)\right)=IH$.\\
			Ta có $AB\parallel A'{B}'$ nên
			\allowdisplaybreaks
			\begin{eqnarray*}
				\text{d}\left(AB,A'C\right)&=&\text{d}\left(AB,(A'B'C)\right)\\
				&=&\text{d}\left(I,(A'B'C)\right)=IH=\dfrac{a\sqrt{15}}{5}.
			\end{eqnarray*}
			Khi đó $\dfrac{1}{I{H^2}}=\dfrac{1}{I{C^2}}+\dfrac{1}{I{I'^2}}\Leftrightarrow I{I}'=a\sqrt{3}$.
		}
		{
			\begin{tikzpicture}[scale=1, font=\footnotesize, line join=round, line cap=round, >=stealth]
				\def\h{3}
				\path
				(0,0) coordinate (A)
				(4,0) coordinate (C)
				(1,-1.5) coordinate (B)
				(0,\h) coordinate (A')
				($(B)+(0,\h)$) coordinate (B')
				($(C)+(0,\h)$) coordinate (C')
				($(A)!1/2!(B)$) coordinate (I)
				($(A')!1/2!(B')$) coordinate (I')
				($(C)!0.5!(I')$) coordinate (H)
				;
				\draw (A)--(B)--(C)--(C')--(A')--(B')--(C') (B)--(B')--(C) (A)--(A') (I)--(I');
				\draw[dashed] (A)--(C)--(I)--(H) (A')--(C)--(I');
				\draw
				pic[draw, angle radius=2mm]{right angle=I--H--C}
				pic[draw, angle radius=2mm]{right angle=B--I--C}
				;
				\foreach \x/\g in
				{A/180,B/200,C/0,A'/180,B'/190,C'/0,I/200,I'/170,H/45}
				\fill[black](\x) circle (1.4pt)
				($(\x)+(\g:3mm)$) node[scale=0.8]{$\x$};
			\end{tikzpicture}
		}
		\noindent
		Thể tích $V$ của khối lăng trụ $ABC.A'B'C'$ là $V=\dfrac{a^2\sqrt{3}}{4}\cdot a\sqrt{3}=\dfrac{3a^3}{4}$.
	}
\end{ex}

\begin{ex}%[TT-SoBacNinh-2022-2023]%[Tran Tony, 12EX6-2023]%[2D4K2-4]
	Trong mặt phẳng $Oxy$, cho số phức $z$ thỏa mãn $|z-1+2i|=3$. Tập hợp các điểm biểu diễn của số phức $w=z(1+i)$ là đường tròn nào dưới đây?
	\choice
	{Đường tròn tâm $I(-3;1)$, bán kính $R=3$}
	{\True Đường tròn tâm $I(3;-1)$, bán kính $R=3\sqrt{2}$}
	{Đường tròn tâm $I(3;-1)$, bán kính $R=3$}
	{Đường tròn tâm $I(-3;1)$, bán kính $R=3\sqrt{2}$}
	\loigiai{
		Giả sử $z=x'+y'i,(x',y'\in\mathbb{R})$, $w=x+yi,(x,y\in\mathbb{R})$.
		\allowdisplaybreaks
		\begin{eqnarray*}
			&&w=z(1+i) \Leftrightarrow x+yi=(x'+y'i)(1+i)\\
			&\Leftrightarrow& x+yi=x'-y'+(x'+y')i\\
			&\Leftrightarrow&\heva{& x=x'-y'\\&y=x'+y'}\Leftrightarrow\heva{&x'=\dfrac{x+y}{2}\\&y'=\dfrac{-x+y}{2}}
		\end{eqnarray*}
		Khi đó
		\allowdisplaybreaks
		\begin{eqnarray*}
			&&|z-1+2i|=3\Leftrightarrow\left|\dfrac{x+y}{2}-1+\left(\dfrac{y-x}{2}+2\right)i\right|=3\\
			&\Leftrightarrow& \left(\dfrac{x+y}{2}-1\right)^2+\left(\dfrac{y-x}{2}+2\right)^2=9\\
			&\Leftrightarrow&\left(x+y-2\right)^2+\left(y-x+4\right)^2=36\\
			&\Leftrightarrow&x^2+y^2-6x+2y-8=0\\
			&\Leftrightarrow&\left(x-3\right)^2+\left(y+1\right)^2=18.
		\end{eqnarray*}
		Vậy tập hợp các điểm biểu diễn số phức $w=z(1+i)$ là đường tròn tâm $I(3;-1)$, bán kính $R=3\sqrt{2}$.
	}
\end{ex}

\begin{ex}%[TT-SoBacNinh-2022-2023]%[Tran Tony, 12EX6-2023]%[2D1G5-2]
	Gọi $S$ là tập hợp tất cả các giá trị nguyên của tham số $ m$ để hàm số $$f(x)=\left|x^4+x^3-5x^2-x+m\right|$$ có bốn điểm cực tiểu $x_1$, $x_2$, $x_3$, $x_4$ thỏa mãn $(x_1^2+1)(x_2^2+1)(x_3^2+1)(x_4^2+1)\geq 68$. Tập $S$ có bao nhiêu tập con?
	\choice
	{$16$}
	{$4$}
	{\True $8$}
	{$32$}
	\loigiai{
		Đặt $g(x)=x^4+x^3-5x^2-x+m\Rightarrow g'(x)=4x^3+3x^2-10x-1$.\\
		$g'(x)=0\Leftrightarrow 4x^3+3x^2-10x-1=0\Leftrightarrow\left[\begin{aligned}
			& x=a\approx-1{,}96\\ 
			& x=b\approx-0{,}098\\ 
			& x=c\approx 1{,}31.\\ 
		\end{aligned}\right.$\\
		Ta có bảng biến thiên
		\begin{center}
			\begin{tikzpicture}
				\tkzTabInit[nocadre=false,lgt=1.2,espcl=2.5,deltacl=0.6]
				{$x$/0.7,$g’(x)$/0.7,$g(x)$/2.2}
				{$-\infty$,$a$,$b$,$c$,$+\infty$}
				\tkzTabLine{ ,-,z,+,z,-,z,+, }
				\tkzTabVar{+/$+\infty$,-/$g(a)$,+/$g(b)$,-/$g(c)$,+/$+\infty$}
			\end{tikzpicture}
		\end{center}
		Trong đó $g(a)\approx-10{,}2+m$, $g(b)\approx 0{,}049+m$, $g(c)\approx-4{,}69+m$.\\
		Để hàm số $f(x)=\left|x^4+x^3-5x^2-x+m\right|$ có $4$ điểm cực tiểu điều kiện là
		$$\left\{\begin{aligned}
			& m+0{,}049>0\\ 
			& m-4{,}69<0\\ 
		\end{aligned}\right.\Leftrightarrow-0{,}049<m<4{,}69.$$
		Ta có $x^4+x^3-5x^2-x+m=0\Leftrightarrow{\left(x^2+1\right)^2}+\left(x^2+1\right)\left(x-7\right)-2x+m+6=0$.\\
		Đặt $t=x^2+1,\left(t\ge 1\right).$ ta được phương trình
		\allowdisplaybreaks
		\begin{eqnarray*}
			&&t^2+t\left(x-7\right)-2x+m+6=0\\
			&\Leftrightarrow&t^2+x\left(t-2\right)-7t+m+6=0\\
			&\Leftrightarrow&t^2-7t+m+6=x(t-2)\\
			&\Rightarrow&(t^2-7t+m+6)^2=(t-1)(t-2)^2\\
			&\Leftrightarrow&t^4-15t^3+2(m+33)t^2-2(7m+46)t+m^2+12m+40.
		\end{eqnarray*}
		Bốn điểm cực tiểu $x_1$, $x_2$, $x_3$, $x_4$ thỏa mãn
		\allowdisplaybreaks
		\begin{eqnarray*}
			&&(x_1^2+1)(x_2^2+1)(x_3^2+1)(x_4^2+1)\geq 68\\
			&\Leftrightarrow&t_1t_2t_3t_4\ge 68\\
			&\Leftrightarrow&m^2+12m+40\geq 68\\
			&\Leftrightarrow&\hoac{&m\geq 2\\&m\leq -14.}
		\end{eqnarray*}
		Do $\left\{\begin{aligned}
			&-0{,}049<m<4{,}69\\ 
			& m\in\mathbb{Z}\\ 
			&\left[\begin{aligned}
				& m\ge 2\\ 
				& m\le-14\\ 
			\end{aligned}\right.\\ 
		\end{aligned}\right.$ nên $m=2,m=3,m=4$.\\
		Vậy số tập con của $S$ là $2^3=8$.
	}
\end{ex}

\begin{ex}%[TT-SoBacNinh-2022-2023]%[Tran Tony, 12EX6-2023]%[2H3G2-7]
	Trong KG $Oxyz$, cho mặt cầu $(S)\colon (x-1)^2+(y-1)^2+(z-1)^2=12$ và mặt phẳng $(\alpha)\colon x-2y+2z+11=0$. Lấy điểm $M$ tùy ý trên $(\alpha)$. Từ $M$ kẻ các tiếp tuyến $MA$, $MB$, $MC$ đến mặt cầu $(S)$, với $A$, $B$, $C$ là các tiếp điểm đôi một phân biệt. Khi $M$ thay đổi thì mặt phẳng $(ABC)$ luôn đi qua điểm cố định $H(a;b;c)$. Tổng $a+b+c$ bằng
	\choice
	{$0$}
	{$\dfrac{7}{2}$}
	{$-\dfrac{3}{4}$}
	{\True $2$}
	\loigiai{
	\immini{
	$(S)$ có tâm $I(1;1;1)$, bán kính $R=\sqrt{12}$.\\
	Gọi $M\left(x_0;{y_0};{z_0}\right)$, do $M\in(\alpha)$ nên $$x_0-2y_0+2z_0+11=0\Leftrightarrow x_0=2y_0-2z_0-11.\qquad (1)$$
	Gọi $A(x;y;z)$, ta có tam giác $IAM$ vuông tại $A$ nên có
	\allowdisplaybreaks
	\begin{eqnarray*}
	&&MA^2=IM^2-IA^2\\
	&\Leftrightarrow&(x-x_0)^2+(y-y_0)^2+(z-z_0)^2=(x_0-1)^2+(y_0-1)^2+(z_0-1)^2-12\\
	&\Leftrightarrow&x^2+y^2+z^2-2x_0x-2y_0y-2z_0z+2x_0+2y_0+2z_0+9=0\qquad (2)
	\end{eqnarray*}
		}{
			\begin{tikzpicture}[scale=0.8,line join=round, line cap=round,thick]
				\def\a{2}
				\def\b{0.5}
				\coordinate (B) at (0,0);
				\coordinate (A) at ($(B)+(2*\a,0)$);
				\coordinate (H) at ($(A)!0.5!(B)$);
				\coordinate (I) at ($(H)+(0,1.3)$);
				\draw[dashed,thin] (A) arc (0:180:\a cm and \b cm);
				\draw (A) arc (0:-180:\a cm and \b cm);
				\path (A) arc (0:-120:{\a} and {\b})	coordinate (C)
				;
				\draw
				let	\p1=($(A)-(I)$), \n1={veclen(\x1,\y1)}, \n2={atan2(\y1,\x1)}
				in (A) arc (\n2:{\n2+240}:\n1);
				\draw[dashed]
				let	\p1=($(A)-(I)$), \n1={veclen(\x1,\y1)}, \n2={atan2(\y1,\x1)}
				in (A) arc (\n2:{\n2-120}:\n1);
				\coordinate (T1) at ($(A)!1cm!90:(I)$);
				\coordinate (T2) at ($(A)!0cm!-90:(I)$);
				\coordinate (M) at (intersection of I--H and A--T1);
				\draw[dashed,thin](H)--(A)--(I)--(M);
				\draw (M)--(A) (M)--(B) (M)--(C);
				\foreach \i/\g in {A/-10,B/200,C/90,H/180,I/90,M/-90}{\draw[fill](\i) circle (1.3pt) ($(\i)+(\g:3mm)$) node[scale=0.7]{$\i$};}
			\end{tikzpicture}
		}
		\noindent
		Mặt khác, $A\in (S)$ nên có $$(x-1)^2+(y-1)^2+(z-1)^2=12\Leftrightarrow x^2+y^2+z^2-2x-2y-2z-9=0.\qquad (3)$$
		Tọa độ điểm $A$ thỏa mãn phương trình $(2)$ và $(3)$ nên trừ vế với vế của hai phương trình đó ta được
		$$(x_0-1)x+(y_0-1)y+(z_0-1)z-x_0-y_0-z_0-9=0.\qquad (4)$$
		Thế $(1)$ vào $(4)$ ta được
		$(2y_0-2z_0-12)x+(y_0-1)y+(z_0-1)z-3y_0+z_0+2=0.\qquad (5)$\\
		Suy ra mặt phẳng $(ABC)$ có phương trình là phương trình $(5)$.\\
		Mặt phẳng $(ABC)$ luôn đi qua điểm cố định $H(a;b;c)$ nên có
		\allowdisplaybreaks
		\begin{eqnarray*}
			&&(2y_0-2z_0-12)a+(y_0-1)b+(z_0-1)c-3y_0+z_0+2=0,\forall \,y_0,z_0\\
			&\Leftrightarrow&(2a+b-3)y_0+(-2a+c+1)z_0-12a-b-c+2=0,\forall \,y_0,z_0\\
			&\Leftrightarrow&\heva{&2a+b-3=0\\&-2a+c+1=0\\&-12a-b-c+2=0}\Leftrightarrow\heva{&a=0\\&b=3\\&c=-1}\Rightarrow H(0;3;-1).
		\end{eqnarray*}
		Vậy $a+b+c=2$.
	}
\end{ex}

\begin{ex}%[TT-SoBacNinh-2022-2023]%[Tran Tony, 12EX6-2023]%[2D3G3-3]
	Cho hàm số $f(x)$ có đạo hàm liên tục trên $\mathbb{R}$ và thỏa mãn các điều kiện $f(0)=-2$, $\left(x^2+1\right){f}'(x)+xf(x)=-x$, $\forall x\in\mathbb{R}$. Gọi $(H)$ là hình phẳng giới hạn bởi đồ thị hàm số $g(x)=\dfrac{1}{1+f(x)}$, hai trục toạ độ và đường thẳng $ x=3$. Quay $(H)$ quanh trục $Ox$ ta được khối tròn xoay có thể tích bằng $V$ (đơn vị thể tích). Khẳng định nào sau đây đúng?
	\choice
	{$ V\in\left(5;9\right)$}
	{$ V\in\left(15;20\right)$}
	{$ V\in\left(11;13\right)$}
	{\True $ V\in\left(35;38\right)$}
	\loigiai{
		Hàm số $f(x)$ có đạo hàm liên tục trên $\mathbb{R}$, ta có
		\allowdisplaybreaks
		\begin{eqnarray*}
			&&\left(x^2+1\right){f}'(x)+xf(x)=-x\\
			&\Leftrightarrow&\sqrt{x^2+1}{f}'(x)+\dfrac{x}{\sqrt{x^2+1}}f(x)=-\dfrac{x}{x^2+1}\\
			&\Leftrightarrow&\left(\sqrt{x^2+1}f(x)\right)'=\left(-\sqrt{x^2+1}\right)'\\
			&\Leftrightarrow&\sqrt{x^2+1}f(x)=-\sqrt{x^2+1}+C.
		\end{eqnarray*}
		Vì $f(0)=-2$ nên $C=-1$ . Suy ra $f(x)=-1-\dfrac{1}{\sqrt{x^2+1}}$.\\
		Khi đó $g(x)=\dfrac{1}{1+f(x)}=-\sqrt{x^2+1}$.\\
		Thể tích khối tròn xoay là $V=\pi\displaystyle\int\limits_0^3g^2(x)\mathrm{\,d}x=\pi\displaystyle\int\limits_0^3\left(-\sqrt{x^2+1}\right)^2\mathrm{\,d}x=12\pi $.\\
		Vậy $V=12\pi\approx 37{,}699\in\left(35;38\right)$.
	}
\end{ex}

\begin{ex}%[TT-SoBacNinh-2022-2023]%[Tran Tony, 12EX6-2023]%[2D4G5-2]
	Xét các số phức $z$ và $w$ thỏa mãn $\left|z\right|=\left|w\right|=1$, $\left|z+w\right|=\sqrt{2}$. Giá trị nhỏ nhất của biểu thức $P=\left|w-\dfrac{4}{z}+2\left(1+\dfrac{w}{z}\right)i\right|$ thuộc khoảng nào?
	\choice
	{\True $\left(2;3\right)$}
	{$\left(4;5\right)$}
	{$\left(3;4\right)$}
	{$\left(7;8\right)$}
	\loigiai{
		Gọi $A$, $B$ lần lượt là biểu diễn của $z$ và $w$ trên mặt phẳng phức $\Rightarrow OA=OB=1$\\
		Ta có $\left|z+w\right|^2+\left|z-w\right|^2=2\left(\left|z\right|^2+\left|w\right|^2\right)\Leftrightarrow 2+\left|z-w\right|^2=4\Rightarrow\left|z-w\right|=\sqrt{2}\Rightarrow AB=\sqrt{2}$.\\
		Nên tam giác $OAB$ vuông cân tại $O$.\\
		Đặt $z=\cos x+i\sin x$ và $w=\cos\left(x-\dfrac{\pi}{2}\right)+\sin\left(x-\dfrac{\pi}{2}\right)i$ hay $w=\sin x+\cos xi\Rightarrow A\left(\cos x;\sin x\right)$ và $B\left(\sin x;\cos x\right)$\\
		$P=\left| w-\dfrac{4}{z}+2\left(1+\dfrac{w}{z}\right)i\right|=\dfrac{\left| wz-4+2iz+2wi\right|}{\left| z\right|}=\left| wz+2iz-4+2wi\right|=\left|\left(w+2i\right)\left(z+2i\right)\right|$\\
		$P=\left|\left(w+2i\right)\left(z+2i\right)\right|=\left| w+2i\right|\left| z+2i\right|=CA\cdot CB$ với $C\left(0;-2\right)$.\\
		Ta có $\heva{&\overrightarrow{CA}=\left(\cos x;\sin x-2\right)\\&\overrightarrow{CB}=\left(\sin x;\cos x-2\right)}\Rightarrow\heva{&CA=\sqrt{\cos^2x+\left(\sin x+2\right)^2}=\sqrt{4\sin x+5}\\&CB=\sqrt{\sin^2x+\left(\cos x+2\right)^2}=\sqrt{4\cos x+5}.}$\\
		Khi đó $P=\sqrt{4\sin x+5}\sqrt{4\cos x+5}=\sqrt{16\sin x\cos x+20\left(\sin x+\cos x\right)+25}$.\\
		Đặt $t=\sin x+\cos x\Rightarrow\heva{&	t\in\left[-\sqrt{2};\sqrt{2}\right]\\&2\sin x\cos x=t^2-1.}$\\
		Khi đó $P=\sqrt{8t^2+20t+17}=\sqrt{8\left(t+\dfrac{5}{4}\right)^2+\dfrac{9}{2}}\ge\dfrac{3\sqrt{2}}{2}$.\\
		Đẳng thức xảy ra khi $\sin x+\cos x=\dfrac{-5}{4}$.\\
		Vậy giá trị nhỏ nhất của $P$ bằng $\dfrac{3\sqrt{2}}{2}\approx 2{,}12$.
	}
\end{ex}

\begin{ex}%[TT-SoBacNinh-2022-2023]%[Tran Tony, 12EX6-2023]%[2D2G6-5]
	Có bao nhiêu cặp số nguyên $(x;y)$ thỏa mãn $$\log_2\dfrac{1+x^2+y^2}{x-2y}\le{4^{x-2y}}-2\cdot{2^{x^2+y^2}}+1?$$
	\choice
	{$6$}
	{\True $13$}
	{$21$}
	{$9$}
	\loigiai{
		Điều kiện $x-2y>0.\qquad (1)$
		\allowdisplaybreaks
		\begin{eqnarray*}
			&&\log_2\dfrac{1+x^2+y^2}{x-2y}\le{4^{x-2y}}-2.2^{x^2+y^2}+1\\
			&\Leftrightarrow&\log_2(1+x^2+y^2)-\log_2(x-2y)\le{2^{2x-4y}}-2^{1+x^2+y^2}+1\\
			&\Leftrightarrow&\log_2\left(1+x^2+y^2\right)+2^{1+x^2+y^2}\le{\log_2}\left(2x-4y\right)+2^{2x-4y}.\qquad (*)
		\end{eqnarray*}
		Xét hàm số $f(t)=\log_2t+2^t,t>0$.\\
		Ta có $f'(t)=\dfrac{1}{t\ln 2}+2^t\ln 2>0,\forall t>0$ nên hàm số $f(t)$ đồng biến trên $\left(0;+\infty\right)$.\\
		Khi đó $(*)\Leftrightarrow 1+x^2+y^2\le 2x-4y\Leftrightarrow{\left(x-1\right)^2}+\left(y+2\right)^2\le 4.\qquad (**)$\\
		Có $\left(x-1\right)^2\le 4\Leftrightarrow-2\le x-1\le 2\Leftrightarrow-1\le x\le 3$. Mà $x\in\mathbb{Z}$ nên $x\in\left\{-1;0;1;2;3\right\}$.
		\begin{itemize}
			\item Với $x=-1\Rightarrow\left(**\right)\Leftrightarrow{\left(y+2\right)^2}\le 0\Leftrightarrow y=-2$ (thỏa mãn $(1)$).
			\item Với $x=0\Rightarrow\left(**\right)\Leftrightarrow{\left(y+2\right)^2}\le 3,y\in\mathbb{Z}\Rightarrow y\in\left\{-3;-2;-1\right\}$ (thỏa mãn $(1)$).
			\item Với $x=1\Rightarrow\left(**\right)\Leftrightarrow{\left(y+2\right)^2}\le 4,y\in\mathbb{Z}\Rightarrow y\in\left\{-4;-3;-2;-1;0\right\}$ (thỏa mãn $(1)$).
			\item Với $x=2\Rightarrow\left(**\right)\Leftrightarrow{\left(y+2\right)^2}\le 3,y\in\mathbb{Z}\Rightarrow y\in\left\{-3;-2;-1\right\}$ (thỏa mãn $(1)$).
			\item Với $x=3\Rightarrow\left(**\right)\Leftrightarrow{\left(y+2\right)^2}\le 0\Leftrightarrow y=-2$ (thỏa mãn $(1)$).
		\end{itemize}
		Vậy có tất cả $13$ cặp số nguyên $\left(x;y\right)$ thỏa đề bài.\\
		\textit{Chú ý}: có thể thay việc đối chiếu điều kiện $(1)$ ở tất cả các trường hợp trên bằng nhận xét: Hình tròn tâm $I(1;-2)$, bán kính $R=2$ nằm hoàn toàn bên dưới đường thẳng $x-2y=0$, nên tất cả các điểm $M(x;y)$ thuộc hình tròn đều có tọa độ thỏa $(1)$.
	}
\end{ex}

\begin{ex}%[TT-SoBacNinh-2022-2023]%[Tran Tony, 12EX6-2023]%[2D2G6-5]
	Gọi $S$ là tập tất cả các giá trị nguyên của tham số $m\in\left[-100;100\right]$ sao cho bất phương trình $3^{x^2-2x+1}-\log_5\left(x^2-2x+6\right)^8+10-\sqrt{-x^2+2x+m}<0$ có nghiệm thực. Tổng tất các các phần tử của $S$ bằng
	\choice
	{$5044$}
	{$5022$}
	{$4914$}
	{\True $5014$}
	\loigiai{
	Ta có
		\allowdisplaybreaks
		\begin{eqnarray*}
			&&3^{x^2-2x+1}-\log_5\left(x^2-2x+6\right)^8+10-\sqrt{-x^2+2x+m}<0\\
			&\Leftrightarrow&3^{\left(x-1\right)^2}-8\log_5\left[\left(x-1\right)^2+5\right]+10-\sqrt{-\left(x-1\right)^2+1+m}<0.
		\end{eqnarray*}
		Đặt $t=\left(x-1\right)^2$, $t\ge 0$. Bất phương trình trở thành $3^t-8\log_5\left(t+5\right)+10-\sqrt{-t+1+m}<0$.\\
		Xét hàm số $f(t)=3^t-8\log_5\left(t+5\right)+10-\sqrt{-t+1+m}$, với $t\ge 0$ .\\
		Ta có $f'(t)=3^t\ln 3-\dfrac{8}{\left(t+5\right)\ln 5}+\dfrac{1}{2\sqrt{-t+1+m}}$.\\
		Với mọi $t\ge 0$ ta có $t+5\ge 5\Leftrightarrow\left(t+5\right)\ln 5>8\Rightarrow \dfrac{8}{\left(t+5\right)\ln 5}<1$.\\
		Suy ra $3^t\ln 3-\dfrac{8}{\left(t+5\right)\ln 5}>0,\,\forall t\ge 0\Rightarrow f'(t)=3^t\ln 3-8.\dfrac{1}{\left(t+5\right)\ln 5}+\dfrac{1}{2\sqrt{-t+1+m}}>0,\,\forall t\ge 0$.\\
		Do đó hàm số $y=f(t)$ đồng biến trên $\left[0;+\infty\right)$ $\Rightarrow f(t)\ge f(0)=3-\sqrt{1+m}$, $\forall t\ge 0$.\\
		Do đó bất phương trình có nghiệm thực khi và chỉ khi $3-\sqrt{1+m}<0\Leftrightarrow 1+m>9\Leftrightarrow m>8$.\\
		Mà $m\in\mathbb{Z},\,m\in\left[-100;\,100\right]$ $\Rightarrow m\in\left\{ 9;\,10;\,11;\,\ldots;\,100\right\}$.\\
		Tổng các phần tử của tập $S$ là $9+10+11+\cdots+100=\dfrac{\left(9+100\right)92}{2}=5014$.
	}
\end{ex}
\Closesolutionfile{ans}
\begin{indapan}{11}
	{ans/ans-2-TT-2-So-BacNinh-23}
\end{indapan}
